\chapter{Future Work}

\section{Purpose of this Chapter}

The AIM framework presented in this manuscript establishes a
mathematical and conceptual foundation for alignment-centred modelling.

However, many aspects remain open and provide opportunities for further
research, development, and practical implementation.

This chapter outlines potential directions for future work, ranging from
mathematical extensions to software architecture and engineering
applications.


\section{Extension of the Dynamic Model}

\subsection{From simplified dynamics to vehicle models}

The current formulation relies on simplified dynamic quantities such as
effective lateral acceleration and jerk.

Future work should extend this toward more realistic vehicle models,
including:
\begin{itemize}
\item multi-body vehicle dynamics,
\item wheel--rail interaction,
\item suspension behaviour,
\item track elasticity.
\end{itemize}

\subsection{Speed-dependent modelling}

The integration of variable speed profiles
\[
v(s)
\]
opens the possibility for more realistic evaluation and optimization of
alignments under operational conditions.


\section{Schwerpunkt-Trassierung}

\subsection{Conceptual development}

The idea of designing the trajectory of a dynamically relevant point,
such as the center of mass, represents a fundamental extension of
classical alignment design.

Further work is required to:
\begin{itemize}
\item formalize the mapping from track geometry to vehicle trajectory,
\item incorporate full dynamic models,
\item define appropriate objective functions.
\end{itemize}

\subsection{Integration into optimization}

A key research direction is the coupling of Schwerpunkt-Trassierung with
optimization methods, enabling:
\begin{itemize}
\item comfort-driven design,
\item dynamic optimization,
\item direct evaluation of vehicle response.
\end{itemize}


\section{Advanced Optimization Methods}

\subsection{SQP extensions}

While Sequential Quadratic Programming provides a powerful baseline,
future work may include:
\begin{itemize}
\item tailored SQP formulations for sparse alignment structures,
\item improved constraint handling,
\item adaptive weighting strategies.
\end{itemize}

\subsection{Global optimization}

Alignment optimization problems are often non-convex.

Future research may explore:
\begin{itemize}
\item global optimization techniques,
\item multi-start strategies,
\item hybrid deterministic--stochastic methods.
\end{itemize}


\section{Network-Level Modelling}

\subsection{From single alignment to network}

Real railway systems consist of interconnected alignments forming
complex networks.

Future work should extend the AIM framework to:
\begin{itemize}
\item branching structures,
\item switches and crossings,
\item corridor-wide optimization.
\end{itemize}

\subsection{Topological constraints}

A systematic treatment of topological constraints remains a major
challenge.

This includes:
\begin{itemize}
\item node consistency,
\item multi-track coupling,
\item network-wide feasibility conditions.
\end{itemize}


\section{Constraint and Objective Libraries}

\subsection{Constraint formalization}

Although many railway rules can be expressed as constraints, a complete
formalization of regulatory frameworks remains an open task.

Future work should aim at:
\begin{itemize}
\item systematic encoding of standards,
\item validation against real projects,
\item parameterization for different railway systems.
\end{itemize}

\subsection{Objective design}

The design of objective functions is crucial for optimization-based
alignment design.

Potential extensions include:
\begin{itemize}
\item cost models,
\item maintenance optimization,
\item energy efficiency,
\item lifecycle considerations.
\end{itemize}


\section{Integration with BIM and Data Standards}

\subsection{Deeper BIM integration}

The integration of AIM with BIM workflows can be further developed by:
\begin{itemize}
\item tighter coupling with IFC alignment models,
\item bidirectional synchronization,
\item embedding AIM states within BIM objects.
\end{itemize}

\subsection{Standard extensions}

The AIM framework may also inspire extensions to existing standards,
for example:
\begin{itemize}
\item inclusion of dynamic properties,
\item support for optimization metadata,
\item representation of control-based alignment models.
\end{itemize}


\section{Interactive and Real-Time Systems}

\subsection{Live optimization}

One of the most promising directions is the development of systems that
allow real-time feedback during alignment design.

This includes:
\begin{itemize}
\item live preview during optimization,
\item interactive constraint adjustment,
\item immediate visualization of dynamic effects.
\end{itemize}

\subsection{User interaction}

The grabbeltisch concept may be extended to:
\begin{itemize}
\item interactive constraint definition,
\item visual debugging of alignment states,
\item intuitive editing of transition parameters.
\end{itemize}


\section{AI and Assisted Design}

\subsection{Data-driven approaches}

Machine learning techniques may complement the AIM framework by:
\begin{itemize}
\item learning from existing alignments,
\item suggesting initial solutions,
\item identifying patterns in design choices.
\end{itemize}

\subsection{Hybrid systems}

A promising direction is the combination of:
\begin{itemize}
\item physics-based modelling (AIM),
\item data-driven methods (AI),
\item interactive user control.
\end{itemize}


\section{AXTRAN Revival and Beyond}

Classical systems such as AXTRAN demonstrated the feasibility of
optimization-based alignment design.

Future work may aim to:
\begin{itemize}
\item reconstruct and generalize such systems,
\item extend them to network-level problems,
\item integrate modern numerical and computational techniques.
\end{itemize}


\section{Software Architecture and Deployment}

\subsection{Scalability}

Practical applications require scalable implementations capable of
handling large datasets and complex networks.

\subsection{Multi-view systems}

Future systems may include:
\begin{itemize}
\item synchronized multi-window views,
\item combined 2D/3D representations,
\item integration of GIS and CAD functionality.
\end{itemize}


\section{Long-Term Vision}

In the long term, the AIM framework may contribute to a new paradigm of
infrastructure modelling in which:
\begin{itemize}
\item alignment becomes the central organizing structure,
\item geometry, dynamics, and optimization are unified,
\item engineering design becomes an interactive computational process.
\end{itemize}


\section{Summary}

\begin{summary}
The AIM framework opens a wide range of research and development
directions.

These include extensions of the dynamic model, integration with BIM,
network-level optimization, real-time systems, and AI-assisted design.

Together, these directions point toward a future in which railway
alignment is no longer treated as static geometry, but as a dynamic,
computational, and optimizable system at the core of infrastructure
engineering.
\end{summary}
